% from ega4-17-fr.tex
\oldpage[IV-4]{64}
$Y$-morphisme $h=g'\circ s':Y'\to X$ (où $g':X'\to X$ est la projection canonique), de sorte que $s'=(1_{Y'},h)_X$, et alors le diagramme
\[
\tag{17.4.1.1}
\label{IV.17.4.1.1}
\xymatrix{
Y' \ar[r]^{s'} \ar[d]_h & X'=Y'\times_Y X \ar[d]^{h\times 1_X}\\
X \ar[r]_{\Delta_f} & X\times_Y X
}
\]
identifie $Y'$ au produit des $(X\times_Y X)$-préschémas $X$ et $X'$. Par suite (I, 4.3.2) si $\Delta_f$ est un isomorphisme local au point $x$, $s'$ est un isomorphisme local au point $y'$ (puisque $x=h(y')$), ce qui prouve que b) implique b''). La réciproque s'obtient en appliquant b'') au cas où l'on prend $Y'=X$, $y'=x$, $g=f$ et $s'=\Delta_f$.

Pour achever la démonstration de (17.4.1), il suffit de prouver les implications
\[
d'')\Rightarrow c)\Rightarrow b)\Rightarrow d'').
\]

$d'')\Rightarrow c)$: Comme $\Omega^1_{X/Y}$ est un $\sh{O}_X$-Module de type fini, il résulte du lemme de Nakayama que la condition c) équivaut à $(\Omega^1_{X/Y})_x/\mathfrak{m}_y(\Omega^1_{X/Y})_x=0$, c'est-à-dire (16.4.5) à $(\Omega^1_{X_y/\Spec(k(y))})_x=0$. On est donc ramené au cas où $Y$ est le spectre d'un corps $k$ et $X$ un $k$-préschéma algébrique. L'hypothèse que $\sh{O}_{X,x}$ est un corps $k'$, extension finie de $k$, entraîne d'abord que $x$ est \emph{fermé} dans $X$ (I, 6.4.2), puis que $x$ est point \emph{maximal} du préschéma noethérien $X$, donc est un point isolé de $X$. Remplaçant $X$ par l'ouvert $\{x\}$ de $X$, on peut donc supposer que $X=\Spec(k')$; mais alors l'hypothèse que $k'$ est extension \emph{finie séparable} de $k$ entraîne $\Omega^1_{k'/k}=0$ (0, 20.6.20), ce qui prouve c).

$c)\Rightarrow b)$: On a vu plus haut que l'on a alors $\Omega^1_{X/Y}|U=0$ pour un voisinage ouvert $U$ de $x$ dans $X$; l'assertion b) résulte alors de la définition de $\Omega^1_{X/Y}$ (16.3.1) et de (16.1.9).

$b)\Rightarrow d'')$: Remplaçant $X$ par un voisinage ouvert de $x$, on peut supposer que $\Delta_f$ est une immersion ouverte; si l'on désigne par $f_y:X_y\to\Spec(k(y))$ le morphisme déduit de $f$ par changement de base, $\Delta_{f_y}$ est alors aussi une immersion ouverte (I, 5.3.4), et comme la condition d'') ne concerne que le préschéma $X_y$, on voit qu'on peut se borner au cas où $Y$ est le spectre d'un corps $k$, $X$ le spectre d'une $k$-algèbre $A$ de type fini; la propriété d'') sera établie si l'on prouve que $A$ est une $k$-algèbre \emph{finie et séparable}, une telle algèbre étant composée directe d'extensions finies séparables de $k$. Si $K$ est une extension algébriquement close de $k$, il revient au même de dire que $A\otimes_k K$ est une $K$-algèbre finie et séparable (4.6.1), donc on voit qu'on peut se borner au cas où $k$ est algébriquement clos. Montrons d'abord que $A$ est une $k$-algèbre \emph{finie}: il suffira de montrer que tout point \emph{fermé} $x$ de $X$ est isolé, car alors l'ensemble de ces points est ouvert dans $X$ et discret, donc fini puisqu'il est quasi-compact ($X$ étant noethérien), ce qui établira notre assertion en vertu de (I, 6.4.4). Or, on a alors $k(x)=k$ puisque $k$ est algébriquement clos (I, 6.4.3), donc il y a une $Y$-section $s$ de $X$ telle que $s(Y)=\{x\}$, et en vertu de (17.4.1.1), $\{x\}$ est l'image réciproque de la diagonale $\Delta_X(X)$ par un morphisme $X\to X\times_Y X$, donc $\{x\}$ est ouvert dans $X$ en vertu de l'hypothèse b). On a ainsi montré que $A$ est une

